\section{Methods}
\label{sec:methods}

\subsection{Geometry and Material Properties}

The abdominal cavity is modelled as an ellipsoidal shell with semi-major
axis $a = \SI{0.15}{\metre}$ and semi-minor axis $b = \SI{0.10}{\metre}$,
corresponding to a typical adult abdomen. The wall thickness is
$h = \SI{15}{\milli\metre}$.

The abdominal wall is treated as an isotropic, linear elastic solid with
properties in the ranges summarised in \cref{tab:materials}.

\begin{table}[htbp]
  \centering
  \caption{Material property ranges for the abdominal wall and cavity fluid.}
  \label{tab:materials}
  \begin{tabular}{@{}lll@{}}
    \toprule
    Property & Range & Source \\
    \midrule
    Young's modulus $E$ & \SIrange{20}{100}{\kilo\pascal} & \citet{Fung1993} \\
    Poisson's ratio $\nu$ & 0.45--0.49 & Soft tissue \\
    Density $\rho_s$ & \SIrange{1000}{1100}{\kilogram\per\cubic\metre} & Tissue \\
    Sound speed $c$ & \SI{1500}{\metre\per\second} & Fluid \\
    Fluid density $\rho_f$ & \SI{1000}{\kilogram\per\cubic\metre} & Water \\
    \bottomrule
  \end{tabular}
\end{table}

\subsection{Mesh Generation}

Three-dimensional finite element meshes are generated using Gmsh
\citep{Geuzaine2009}. The geometry consists of two concentric ellipsoids
(inner surface for the fluid domain, outer surface for the structural domain).

% TODO: Mesh details, element type, mesh convergence results

\subsection{Structural Modal Analysis}

The structural eigenvalue problem is:
%
\begin{equation}
  \mathbf{K} \boldsymbol{\phi} = \omega^2 \mathbf{M} \boldsymbol{\phi},
  \label{eq:structural_eigenvalue}
\end{equation}
%
where $\mathbf{K}$ and $\mathbf{M}$ are the assembled stiffness and mass
matrices. This is solved using FEniCSx \citep{Baratta2023} with the
Krylov--Schur algorithm in SLEPc \citep{Hernandez2005}.

% TODO: Function space, element type, boundary conditions, solver settings

\subsection{Acoustic Modal Analysis}

% TODO: Helmholtz equation discretisation, rigid-wall BCs

\subsection{Coupled Fluid--Structure Analysis}

% TODO: Mixed formulation, coupling conditions

\subsection{Mesh Convergence Study}

% TODO: Richardson extrapolation, GCI results
